% =====================================================================
%  PART I -- ORIENTATION
%
%  Preface, motivation, intuition, source clarification, literature graph.
%  No symbols are required to read this part.  Nothing in it is assumed by
%  Part II; it exists so that Part II is read for the right reasons.
% =====================================================================

\chapter{The Claim}\label{ch:claim}

\section{The founding hypothesis}\label{sec:hypothesis}

Cognitive Space Theory rests on one hypothesis, and the whole apparatus of
this book is a consequence of taking it seriously.

\begin{quote}\itshape
Meaning admits no global chart.
\end{quote}

That is: there is no single coordinate system in which everything that can be
meant is simultaneously and consistently expressible. Meaning is
coordinatizable only \emph{locally}, in patches, and the patches do not agree
on their overlaps.

\begin{intu}\Intu
A map of the Earth is a sheet of paper with coordinates on it. No single
sheet covers the globe without tearing or distortion, so cartographers use an
\emph{atlas}: many local sheets, plus rules for translating a position on one
sheet into a position on another. The rules are the real content. The sheets
are interchangeable; the translations are where the curvature of the Earth
shows up.

The hypothesis above says that meaning is like the globe and not like the
sheet. If it were like the sheet --- if one big coordinate system sufficed ---
then a cognitive state would be a vector, cognition would be linear algebra in
a fixed basis, and this book would have nothing to say. Everything here exists
because the translations between local frames fail to compose.
\end{intu}

The hypothesis is not decoration and it is not unfalsifiable. It has a precise
mathematical shadow, developed in Part~\ref{part:foundations}: if meaning
admitted a global chart, then the complex would have one vertex, the sheaf
would have no restriction maps, the coboundary would be zero, the Laplacian
would be zero, and every state would be coherent. The theory would collapse to
its own trivial case. \emph{Structure exists exactly to the extent that
local coordinatizations fail to agree.} Chapter~\ref{ch:complex},
Example~\ref{ex:atom}, computes that collapse explicitly, so that the reader
can see the degenerate case and know what the non-degenerate one is departing
from.

\begin{hon}\Hon
\Proposed\ The founding hypothesis is a modelling commitment, not a theorem,
and nothing in this book proves it. What the book does is make it
\emph{consequential}: every construction is derived from it, so that if the
hypothesis is wrong, the constructions are not merely unmotivated but
provably vacuous. That is the strongest honest status a founding hypothesis
can have.
\end{hon}

\section{What is being modelled}

Three words are used throughout and each is used in a restricted, technical
sense. They are declared here because their ordinary meanings would import
claims the theory does not make.

\begin{center}
\begin{tabular}{@{}>{\raggedright\arraybackslash}p{0.20\textwidth}%
                  >{\raggedright\arraybackslash}p{0.73\textwidth}@{}}
\toprule
\textbf{region} &
  A place where meaning is locally coordinatizable. Formally a vertex of the
  complex. \emph{Not} an agent, a module, a subsystem, or a processor.\\
\addlinespace[3pt]
\textbf{content} &
  What is expressible within one region, in that region's own coordinates.
  Formally an element of a stalk. \emph{Not} a belief; a region does not
  hold an opinion, because a region is not the kind of thing that can have
  one.\\
\addlinespace[3pt]
\textbf{thought} &
  A choice of content in every region that survives every comparison.
  Formally a global section. There is \emph{one} of these at a time, expressed
  in many local coordinate systems --- not one per region.\\
\bottomrule
\end{tabular}
\end{center}

\begin{hon}\Hon
The third row is the one that is easiest to get wrong, and getting it wrong
changes the mathematics rather than merely the prose. If each region holds its
own belief, then disagreement between regions is a normal and permanent
feature to be negotiated, and the natural theory is a theory of consensus.
That theory exists, it is good, and it is \emph{not this one}: see
Hansen--Ghrist's discourse sheaves~\cite{HansenGhristOpinion}, discussed in
\S\ref{sec:neighbours}, where the vertices are genuinely separate agents.

Under the reading taken here there is one thought and many coordinate systems,
so disagreement is not a social fact but a \emph{defect of the atlas}: either
the frames can be adjusted until they agree, or they cannot, and the
obstruction to adjusting them is a topological invariant. That is what makes
the dynamics of Part~\ref{part:foundations} onwards possible at all. A
consensus theory has nothing to converge to except agreement; this theory has
something to converge to \emph{and} a classification of the cases where
convergence is impossible.
\end{hon}

\section{The thesis, and what each clause forbids}

\begin{quote}\itshape
A cognitive space is a cellular sheaf on a finite simplicial $2$-complex.
Its state is a section, which relaxes. Its restriction maps are the theory's
only carrier of learned structure, and they mold. Its complex grows and
prunes, and growth is triggered by an obstruction rather than by a rule.
\end{quote}

Every clause forbids something, and the forbidding is a live constraint on
what may be written later.

\begin{center}
\begin{tabular}{@{}>{\raggedright\arraybackslash}p{0.42\textwidth}%
                  >{\raggedright\arraybackslash}p{0.51\textwidth}@{}}
\toprule
\textbf{clause} & \textbf{forbids} \\
\midrule
sheaf on a \textbf{$2$-complex} &
  dropping to a graph, where the sheaf condition is vacuous
  (Theorem~\ref{thm:functoriality})\\
\addlinespace[2pt]
the state \textbf{relaxes} &
  treating the state as stored data rather than as a dynamical variable\\
\addlinespace[2pt]
restriction maps \textbf{mold} &
  freezing the maps at the identity, which provably empties the theory
  (Proposition~\ref{prop:identityfatal})\\
\addlinespace[2pt]
\textbf{grows and prunes} &
  monotone architectures that only accumulate\\
\addlinespace[2pt]
growth from an \textbf{obstruction} &
  hand-written growth rules; the trigger must be a computed invariant\\
\bottomrule
\end{tabular}
\end{center}

\section{What Cognitive Space Theory is not}\label{sec:isnot}

Stated as refusals, because each of these is a place where an adjacent and
better-established framework would absorb the theory if the boundary were left
implicit.

\begin{enumerate}
  \item \textbf{Not a knowledge graph.} A knowledge graph is a labelled
    relation with no dynamics and no obstruction theory. The working test is
    sharp: \emph{if a proposed feature would behave identically on a labelled
    graph, it does not belong here.}
  \item \textbf{Not a geometric semantics in a single space.} The standard
    geometric account of concepts~\cite{Gardenfors} places them as convex
    regions of one global quality space. Cognitive Space Theory is precisely
    the denial of that global space (\S\ref{sec:hypothesis}). The two are not
    variants; they disagree on whether the global chart exists.
  \item \textbf{Not a compositional semantics.} Categorical accounts of how
    meanings compose~\cite{CoeckeSadrzadehClark} are about \emph{building}
    a meaning from parts. This theory is about meaning \emph{failing} to
    assemble, and about what that failure is a measurement of. A composition
    calculus always composes; the object of study here is the residue when
    assembly does not go through.
  \item \textbf{Not a model of brains.} Mechanisms are occasionally borrowed
    from neuroscience as \emph{design constraints}. No claim is made to
    explain neural data, and nothing here should be evaluated against it.
  \item \textbf{Not a thermodynamic or free-energy theory.} The cost
    functionals introduced later are this theory's own construction. Naming
    them after an established principle would import authority that has not
    been earned.
  \item \textbf{Not a claim about consciousness, agency or life.} Where a
    construction has a suggestive name, the name is a label for a morphism and
    nothing further is asserted.
\end{enumerate}

\begin{agenda}
\agset{
  \item The founding hypothesis is stated in a form whose failure has a
    computable consequence (collapse to Example~\ref{ex:atom}).
  \item The three primitive words are fixed, and the agent reading is
    explicitly excluded with its reason.
  \item The thesis is stated as a list of prohibitions, each pointing at the
    result that enforces it.
}
\agopen{
  \item No independent argument is given \emph{for} the founding hypothesis.
    At present it is motivated by the fact that its negation trivialises the
    theory, which is a consistency argument, not evidence.
  \item ``Region'' is defined by what it is not. A positive criterion --- given
    a domain, what determines how many regions there are and where their
    boundaries fall --- is supplied only operationally, by the ingest
    construction of Chapter~\ref{ch:ingest}, and it inherits that
    construction's arbitrariness in the choice of witness set.
  \item The prohibition ``growth from an obstruction, not a rule'' is stated
    but not yet honoured: the obstruction that would serve as trigger is
    identified only as a candidate in \S\ref{sec:mv}.
}
\agsteps{
  \item Find one phenomenon that a global-chart semantics cannot represent and
    a sheaf semantics can, stated so that the difference is a computation and
    not an interpretation. Contextuality (\S\ref{sec:coefficient}) is the
    leading candidate and would convert O1 from a consistency argument into
    an argument from evidence.
  \item Give a positive criterion for region individuation, or accept
    explicitly that regions are an input to the theory and not an output, and
    say so here rather than leaving it to be inferred.
  \item Close \S\ref{sec:mv} to a theorem; if it closes, restate the fifth
    clause of the thesis as a result rather than a commitment.
}
\end{agenda}

% =====================================================================
\chapter{The Picture Before the Symbols}\label{ch:picture}

Nothing in this chapter is used later. It exists so that the definitions of
Part~\ref{part:foundations} are recognised when they arrive, rather than
merely accepted.

\section{Charts, atlases, and the failure to agree}

\begin{intu}\Intu
Suppose a body of knowledge is covered by overlapping regions, and inside each
region you have a way of writing things down --- a coordinate system private
to that region. Two regions that overlap must be able to compare what they
say about the overlap. Comparison requires a \emph{translation}: a rule taking
what is written in one region's coordinates into a form the overlap can read.

Three things can then happen.
\begin{itemize}
  \item The translations agree everywhere. There is a single consistent thing
    being described, written down in many ways. This is a \emph{thought}.
  \item The translations disagree, but the disagreement can be removed by
    adjusting the translations themselves --- the regions were saying the same
    thing in badly aligned frames. This is a \emph{misalignment}, and it is
    repairable.
  \item The translations disagree, and no adjustment of them removes the
    disagreement, because going around a loop of regions and coming back
    returns you somewhere else. This is a \emph{hole}. It is not repairable by
    translation, only by changing the regions.
\end{itemize}
The entire technical content of this book is: making those three cases
precise, proving that they are exhaustive and mutually orthogonal, computing
which one you are in, and saying what to do in each.
\end{intu}

The three cases are not an informal taxonomy. They are the three summands of
an orthogonal direct-sum decomposition of the space of pairwise comparisons
(Theorem~\ref{thm:hodge}), they are computable by linear algebra, and the
third is a topological invariant that no amount of local adjustment can
change.

\section{A single region cannot be wrong}

Take the smallest possible cognitive space: one region, no overlaps. There are
no comparisons to make, so no comparison can fail, so every content whatsoever
is a thought. The space has no structure and no capacity for error.

This is worth stating early because it is the exact sense in which the theory
is \emph{about} the relations and not about the regions. Content lives in the
regions; \emph{everything else} lives between them. Example~\ref{ex:atom}
carries this out in symbols and finds that every invariant the book defines is
either zero or trivial on the one-region space.

\section{The degenerate corner: what a neural network is, from here}
\label{sec:degenerate}

The framework has a corner in which it reproduces machinery that already
exists, and locating that corner exactly is the most direct way to say what
the rest of the framework is for.

\begin{intu}\Intu
Suppose you make three simplifications at once:
\begin{enumerate}
  \item allow only regions and overlaps of two --- no three-way overlaps;
  \item make every translation the identity, so that all regions share one
    coordinate system after all;
  \item keep only the relaxation dynamics, and never adjust the translations
    or the regions.
\end{enumerate}
What remains is a graph with vectors on its nodes, diffusing toward agreement.
That is a graph neural network, and Proposition~\ref{prop:nnsheaf} makes the
identification exact for the linear case.
\end{intu}

Each simplification has a provable cost, and the three costs are the reason
the rest of the book exists:

\begin{center}
\begin{tabular}{@{}>{\raggedright\arraybackslash}p{0.30\textwidth}%
                  >{\raggedright\arraybackslash}p{0.63\textwidth}@{}}
\toprule
\textbf{simplification} & \textbf{what it provably costs} \\
\midrule
no three-way overlaps &
  the sheaf condition becomes vacuous, so nothing can ever obstruct and growth
  can never be derived (Theorem~\ref{thm:functoriality}); and the three-way
  decomposition loses a summand, so misalignment and holes stop being
  distinguishable (Theorem~\ref{thm:hodge})\\
\addlinespace[3pt]
identity translations &
  the only thought is exact consensus, and relaxation is plain averaging;
  diffusion drives everything to the global mean
  (Proposition~\ref{prop:identityfatal})\\
\addlinespace[3pt]
no adjustment &
  the theory's only carrier of learned structure is frozen, and the sheaf
  becomes a decoration on a graph\\
\bottomrule
\end{tabular}
\end{center}

The second row is not a criticism from outside: it is the published diagnosis
of oversmoothing in graph neural networks, and the published repair is exactly
to make the translations non-trivial~\cite{BodnarSheaf}. The first row has an
independent published counterpart in expressiveness: message passing that sees
higher cells is strictly stronger than message passing that sees only
edges~\cite{BodnarSIN}.

\begin{hon}\Hon
This is the one place where the theory can be checked against an established
literature rather than argued for. It should therefore be the first thing a
sceptical reader is shown, and it is deliberately placed before any
definition. If the identification in Proposition~\ref{prop:nnsheaf} were
wrong, the theory would be in serious trouble; it is not wrong, and it is
proved.
\end{hon}

\section{Reading, as a sequence of operations}\label{sec:reading}

\begin{intu}\Intu
What happens when a body of knowledge absorbs a text? Not storage, and not
appending. The stages, each of which becomes a named operation in
Part~\ref{part:foundations} and later:

\begin{enumerate}
  \item \textbf{Ingest.} The text arrives as a relation --- these things occur
    together in this passage, those in that one --- and a relation determines a
    complex canonically, with no parameters (Chapter~\ref{ch:ingest}).
  \item \textbf{Anchor.} Some of the new material is identified with material
    already present. This is the only step in the whole pipeline containing a
    free choice, and it is where interpretation lives: two readers who anchor
    the same text differently have read different texts.
  \item \textbf{Glue.} The two complexes are joined along what was anchored.
  \item \textbf{Ferment.} The join has properties neither part had. Some
    disagreement is discharged by relaxation, some by adjusting translations
    --- and some cannot be, because the act of joining created a hole that was
    not present in either piece (\S\ref{sec:mv}).
  \item \textbf{Consolidate.} What survives is compressed and retained; what
    does not is pruned.
\end{enumerate}
Step 4 is the interesting one and the reason this is called fermentation
rather than insertion. \emph{New obstructions are a property of the join,
not of either part.} A text you already agreed with entirely teaches you
nothing, and a text you cannot anchor at all teaches you nothing; the content
is in the residue.
\end{intu}

\begin{agenda}
\agset{
  \item The three-case taxonomy (thought / misalignment / hole) is stated
    informally and is discharged by Theorem~\ref{thm:hodge}.
  \item The degenerate corner is located precisely and each simplification is
    paired with the result that prices it.
  \item The reading pipeline is stated as five named stages, four of which are
    given formal definitions later.
}
\agopen{
  \item Stage~5, consolidation, has no formal counterpart in
    Part~\ref{part:foundations}. It is treated in later parts whose
    foundations predate this one.
  \item Stage~2, anchoring, is identified as the locus of interpretation but
    is given no criterion. Nothing here says which anchoring is \emph{better},
    only that the choice exists and matters.
}
\agsteps{
  \item Define anchoring formally as a span with an isomorphism of sheaves on
    its apex (sketched in \S\ref{sec:mv}), and state the invariance question:
    which downstream quantities are independent of the anchoring chosen.
  \item Bring consolidation into Part~\ref{part:foundations} or state
    explicitly that it is deferred, so that stage~5 is not silently missing.
}
\end{agenda}

% =====================================================================
\chapter{Sources, Position, and the Literature Graph}\label{ch:sources}

\section{Naming discipline}\label{sec:naming}

This theory borrows vocabulary from several mature fields. Borrowed words
arrive with expectations attached, and an unearned word is a claim made
silently. The following rules are enforced throughout.

\begin{enumerate}
  \item \textbf{A borrowed term is used only with its source meaning, or not
    at all.} Where a weaker notion is meant, a different word is coined.
  \item \textbf{No term is used whose associated theorems are not being
    invoked.} This retires the word \emph{stratified} from the description of
    the base complex; \S\ref{sec:stratword} explains what was meant, why the
    word was wrong, and where a genuine stratification does appear.
  \item \textbf{Suggestive names are labels, never arguments.} If a
    construction is called reproduction, growth or memory, that name may not
    be used as a premise in any proof.
\end{enumerate}

\begin{hon}\Hon
Terms retired from this book under rule~2, each with its replacement:
\emph{stratified} (of the base complex) $\to$ nothing; the two-scale structure
it named is not assumed, and a genuine stratification appears instead on the
configuration space, \S\ref{sec:stratconf}. \emph{Belief} (of stalk content)
$\to$ \emph{content}, since ``belief'' presupposes a bearer. \emph{Organ},
\emph{agent}, \emph{subsystem} (of vertices) $\to$ \emph{region}, since
vertices are not processors. \emph{Free energy} (of cost functionals) $\to$
the functional's own name.
\end{hon}

\section{The five lineages}

Cognitive Space Theory sits at the meeting of five bodies of work. For each,
what is taken and what is deliberately not taken.

\paragraph{1. Cellular sheaves and their spectral theory.}
Cellular sheaves are due to Shepard~\cite{Shepard} and were developed as a
general applied framework by Curry~\cite{Curry}. The metric theory this book
depends on --- the sheaf Laplacian $\Lap=\cob^{\adjt}\cob$, its spectrum, its
Cheeger-type inequalities --- is Hansen and Ghrist~\cite{HansenGhrist}.
\emph{Taken:} the entire apparatus of Chapters~\ref{ch:sheaf}
and~\ref{ch:laplacian}, essentially unmodified. \emph{Not taken:} their
setting is a fixed complex with fixed stalks. Everything about growth and
molding is outside it and must be built.

\paragraph{2. Data-to-complex constructions.}
Dowker's theorem on the homology of relations~\cite{Dowker}, its functorial
and stable refinement~\cite{ChowdhuryMemoli}, the nerve construction, and the
persistence machinery with its stability
theorems~\cite{CEH,ChazalStructure,Carlsson}. \emph{Taken:} the ingest
operator of Chapter~\ref{ch:ingest}, and the demand that ingest be
\emph{stable}. \emph{Not taken:} persistence itself. This book measures a
complex at one scale, not a filtration across scales; Mapper~\cite{SinghMemoliCarlsson}
is named and rejected, for the reason in \S\ref{sec:alternatives}.

\paragraph{3. Sheaf-theoretic contextuality.}
Abramsky and Brandenburger~\cite{AbramskyBrandenburger} show that
non-locality and contextuality are exactly the non-existence of a global
section of a presheaf of distributions over a cover of measurement contexts;
the cohomological refinement is~\cite{ABKLM}. \emph{Taken:} the identification
of incoherence with a global-section obstruction, and --- decisively --- the
observation that the \emph{coefficient algebra} is a free parameter with
consequences (\S\ref{sec:coefficient}). \emph{Not taken:} any quantum-mechanical
claim. The site is different and no physical interpretation is asserted.

\paragraph{4. Higher-order networks and sheaves on graphs in learning.}
Signal processing beyond pairwise interactions~\cite{Battiston,SchaubTSP};
sheaf neural networks~\cite{HansenGebhart}; neural sheaf
diffusion~\cite{BodnarSheaf}; the expressiveness separation for higher-cell
message passing~\cite{BodnarSIN}. \emph{Taken:} the diagnosis of oversmoothing
as a statement about $\ker\Lap$, which is the sharpest available argument that
non-identity restriction maps are necessary rather than ornamental.
\emph{Not taken:} the learning setting. These are supervised models trained on
labelled data; nothing here is.

\paragraph{5. Synchronization, fusion, and inference of structure.}
The connection Laplacian and vector diffusion maps~\cite{SingerWu}; $O(d)$
synchronization~\cite{BSS}; sheaves for heterogeneous sensor
integration~\cite{RobinsonTSP,RobinsonSensor}; learning sheaf Laplacians from
observed signals~\cite{HansenGhristLearning}. \emph{Taken:} this is the
literature home of \emph{molding} --- adjusting restriction maps toward
consistency is the synchronization problem, and Robinson's consistency radius
is this book's discord in an applied setting. \emph{Not taken:} their
objectives are estimation against a ground truth. There is no ground truth
here.

\section{The literature graph}

\begin{center}
\begin{tikzpicture}[
  scale=1.0,
  node distance=6mm,
  clus/.style={draw,rounded corners=3pt,align=center,font=\footnotesize,
               inner sep=4pt,minimum height=9mm,fill=white},
  core/.style={draw,thick,rounded corners=3pt,align=center,font=\small\bfseries,
               inner sep=5pt,fill=linkblue!10,draw=linkblue},
  take/.style={-{Latex[length=2.2mm]},thick,examplecol},
  drop/.style={-{Latex[length=2.2mm]},dashed,honestycol},
  el/.style={font=\scriptsize,examplecol,inner sep=1pt},
  dl/.style={font=\scriptsize,honestycol,inner sep=1pt}]

  \node[core] (cst) at (0,0) {Cognitive\\Space Theory};

  \node[clus] (dow) at (-4.6,2.6)
     {\textbf{data $\to$ complex}\\ Dowker; Chowdhury--M\'emoli;\\
      nerve; stability \cite{CEH,ChazalStructure}};
  \node[clus] (sh)  at (0,3.3)
     {\textbf{cellular sheaves}\\ Shepard; Curry;\\ Hansen--Ghrist};
  \node[clus] (ctx) at (4.7,2.6)
     {\textbf{contextuality}\\ Abramsky--Brandenburger;\\ ABKLM};
  \node[clus] (syn) at (5.0,-1.4)
     {\textbf{synchronization}\\ \textbf{\& fusion}\\ Singer--Wu; BSS;\\ Robinson};
  \node[clus] (gnn) at (0,-3.0)
     {\textbf{higher-order \& sheaf}\\ \textbf{networks}\\ Bodnar et al.;\\
      Hansen--Gebhart};
  \node[clus] (sem) at (-5.0,-1.4)
     {\textbf{geometric \&}\\ \textbf{compositional semantics}\\
      G\"ardenfors; DisCoCat};
  \node[clus] (lat) at (-2.6,-5.0)
     {\textbf{order-theoretic}\\ \textbf{sheaves}\\ Ghrist--Riess};
  \node[clus] (opn) at (3.2,-5.0)
     {\textbf{discourse sheaves}\\ Hansen--Ghrist};

  \draw[take] (dow) -- node[el,above,sloped] {ingest, stability} (cst);
  \draw[take] (sh)  -- node[el,right] {$\Lap$, $\Hzero$, spectrum} (cst);
  \draw[take] (ctx) -- node[el,above,sloped] {coefficient algebra} (cst);
  \draw[take] (syn) -- node[el,below,sloped] {molding} (cst);
  \draw[take] (gnn) -- node[el,right] {oversmoothing $=\ker\Lap$} (cst);
  \draw[drop] (sem) -- node[dl,above,sloped] {global chart} (cst);
  \draw[drop] (lat) -- node[dl,below,sloped] {no Hodge split} (cst);
  \draw[drop] (opn) -- node[dl,below,sloped] {agent semantics} (cst);

\end{tikzpicture}
\captionof{figure}{Provenance. Solid green: taken, with what is taken on the
edge. Dashed red: a nearest neighbour whose central commitment is
\emph{rejected}, with the point of divergence on the edge. The dashed edges
carry more information than the solid ones, because a theory is located by
what it refuses.}
\label{fig:litgraph}
\end{center}

\section{Nearest neighbours, and exactly where they diverge}
\label{sec:neighbours}

Three bodies of work are close enough that the difference must be stated
precisely rather than gestured at.

\paragraph{Discourse sheaves~\cite{HansenGhristOpinion}.} The closest
published relative: the same Laplacian, the same disagreement functional, on
the same kind of object. The divergence is semantic and it is total. There, a
vertex is an agent with private opinions and the restriction maps express how
an agent's opinion appears in a shared discourse; disagreement is between
people. Here, a vertex is a chart and disagreement is curvature. The
mathematics is shared; the question asked of it is not.

\paragraph{Neural sheaf diffusion~\cite{BodnarSheaf}.} Shares the diagnosis of
identity restriction maps and the remedy of learning non-trivial ones. Two
divergences: their base is a graph, so the sheaf condition never binds
(Theorem~\ref{thm:functoriality}); and their restriction maps are learned by
supervised gradient descent against labels, whereas molding here has no
external target and must be driven by an intrinsic functional.

\paragraph{Sheaf contextuality~\cite{AbramskyBrandenburger}.} Shares the
central identification --- incoherence is the failure of a global section ---
on a different site. The divergence is the coefficient algebra: their stalks
carry distributions over a semiring, in which positivity generates the
obstruction; this book's stalks are vector spaces over a field, chosen so that
a metric and a three-way Hodge decomposition exist. \S\ref{sec:tradeoff}
states what that choice costs, since it is a real cost and not a free
upgrade.

\begin{agenda}
\agset{
  \item Every construction used in Part~\ref{part:foundations} has a named
    lineage, with the take/leave boundary drawn explicitly.
  \item The three nearest neighbours are separated from this theory by a
    stated divergence rather than by tone.
  \item A naming discipline is in force, with the retired terms listed.
}
\agopen{
  \item No lineage is claimed for the dynamics. Molding has a home in the
    synchronization literature, but the two-phase separation of relaxation
    from consolidation does not have one here, and it is not yet established
    whether one exists.
  \item The claim in \S\ref{sec:neighbours} that molding ``must be driven by
    an intrinsic functional'' names no functional. Until one is given and
    shown to have the required stationary points, this is a requirement rather
    than a design.
}
\agsteps{
  \item Search the synchronization literature for an unsupervised objective
    whose stationary points are flat connections, and either adopt it with
    attribution or record that none was found.
  \item Extend Figure~\ref{fig:litgraph} with a sixth cluster for the dynamics
    once O1 is resolved --- or record the absence of one as a finding, since a
    construction with no lineage is either novel or naive and it matters
    which.
}
\end{agenda}
